\section{Conclusiones}

\subsection{Lecciones aprendidas}

\textbf{Problema de negocio resuelto.} MisiOps guarda credenciales y datos financieros de sus usuarios. Una fuga o una dependencia vulnerable destruiría la confianza en el producto. Necesitábamos detectar esos riesgos antes de publicar, no después, y poder desplegar de forma repetible y segura.

\textbf{Limitaciones iniciales.} La seguridad dependía de la revisión manual del PR y de \texttt{pip-audit}. No se revisaban secretos en el historial, ni imágenes, ni configuración, ni la aplicación en ejecución. Los resultados quedaban dispersos y nada impedía integrar un cambio inseguro. Además, la aplicación solo corría en local con Docker Compose y no tenía un entorno en la nube.

\textbf{Qué cambió con DevOps.}
\begin{itemize}
    \item La seguridad corre en cada PR, automatizada y con versiones fijadas de cada herramienta.
    \item Los hallazgos llegan a un solo lugar (DefectDojo), con historial y estado de cada uno.
    \item \texttt{main} solo recibe código que pasó todos los escaneos.
    \item La infraestructura es código probado: un merge a \texttt{develop} actualiza \texttt{dev}, y un merge aprobado a \texttt{main} actualiza \texttt{prod} con exactamente las imágenes que se escanearon.
\end{itemize}

\textbf{Ventajas medidas.}
\begin{itemize}
    \item Se cerraron vulnerabilidades críticas reales (Next.js, certifi, imagen de PostgreSQL): los hallazgos críticos activos bajaron de 39 a 3 y los bloqueantes de 263 a 0.
    \item El pipeline de seguridad completo tarda unos 4 minutos, y el despliegue a \texttt{prod} unos 11 minutos incluyendo la aprobación.
    \item El gate demostró su valor en un caso real: además de bloquear el PR de prueba \#131, bloqueó la promoción a \texttt{main} cuando Debian publicó el fix de dos CVE HIGH en la imagen base del backend.
    \item El rollback automático de ECS actuó en el primer despliegue fallido sin dejar recursos huérfanos.
\end{itemize}

\textbf{Lecciones técnicas.}
\begin{itemize}
    \item Un escáner que no alcanza su objetivo no equivale a ``cero hallazgos'': los jobs fallan explícitamente si ZAP no llega a la aplicación.
    \item Un gate útil bloquea solo lo accionable. Bloquear CVE sin fix publicado frenaría todos los merges sin mejorar la seguridad.
    \item Los títulos de \textit{Test} en DefectDojo deben ser únicos por reporte, porque si no una importación cierra los hallazgos de otra.
    \item Construir las imágenes una sola vez y reutilizarlas para escanear, probar y desplegar garantiza que lo que llega a producción es lo que se analizó.
\end{itemize}

\subsection{Recomendaciones}

\begin{itemize}
    \item Alojar DefectDojo en un servidor propio con credenciales robustas, en lugar del equipo local con \texttt{admin/admin} y un túnel ngrok.
    \item Agregar \textit{rate limiting} y bloqueo temporal en \texttt{/auth/login} (amenazas T1 y T10) y rotación de la \texttt{SECRET\_KEY} (T2).
    \item Al vencer la exención EX-001 (2026-12-31), exigir los checks de seguridad también en \texttt{develop}.
    \item Fijar las GitHub Actions por SHA y activar Dependabot o Renovate para actualizar dependencias de forma continua.
    \item Agregar una bitácora de auditoría de los cambios financieros (amenaza T6).
    \item Restringir la política del rol de ejecución de CloudFormation creado por el \textit{bootstrap} de CDK, que hoy conserva \texttt{AdministratorAccess}.
    \item Migrar la CSP del frontend a una política con \textit{nonce} para eliminar \texttt{'unsafe-inline'}.
\end{itemize}
